\cleardoublepage
\chapter*{Resumen}
\addcontentsline{toc}{chapter}{Resumen}

\textbf{Título:} \Title.\footnote{Trabajo de grado.}

\textbf{Autor:} \Author.\footnote{\FacultyName. \SchoolName. Director: \DirectorName, \DirectorDegree.}

\textbf{Palabras clave:} redes neuronales convolucionales equivariantes, imágenes histopatológicas, procesamiento de imágenes, aprendizaje semi-supervisado.

\textbf{Descripción:}

\begingroup\small

% TODO: reemplazar por un resumen final que sintetice problema, método,
% resultados cuantitativos principales y conclusiones.
La patología digital representa un cambio de paradigma en el enfoque tradicional de diagnóstico médico. Dentro de la medicina contemporánea los avances tecnológicos en curso remodelan constantemente las prácticas de atención médica. Por lo que, el desarrollo de la patología digital es un logro transformador que permite introducir una combinación de dispositivos electrónicos de punta, como los escáneres de láminas histológicas, análisis de datos y experiencia clínica que puede transformar significativamente la forma en que se brinda la atención al paciente.

Dentro de las aplicaciones de la patología digital, una de las que mayor acogida e interés ha provocado han sido las que respectan al cáncer. Esto es debido a que es una de las principales causas de muerte a nivel global. La carga del cáncer aumentará aproximadamente en un 60\% durante las próximas dos décadas, lo que afectará aún más a los sistemas de salud, a las personas y a las comunidades.

Dado esto, se requiere un nuevo enfoque como el uso de modelos de inteligencia artificial (IA), que optimizan y aceleran los métodos actuales. Aunque se han desarrollado proyectos de aprendizaje automático dentro del ámbito de la patología digital, la limitada disponibilidad de imágenes anotadas y segmentadas obstaculiza los procesos de aprendizaje profundo, pues realizar estos procesos de anotación exige gran demanda de tiempo y esfuerzo por parte de los patólogos. Hasta ahora los métodos tradicionales de aprendizaje profundo han solucionado estos problemas con el uso de aumento de datos, que consiste en generación de datos nuevos a partir de los que hay, pero este proyecto pretende mostrar el uso de representaciones equivariantes en los modelos de aprendizaje profundo, que permitirían una representación estable y eficaz de imágenes en patología digital.

Por esto, se propone un modelo de aprendizaje profundo semi-supervisado de redes neuronales equivariantes que permita optimizar el tiempo y los recursos de preprocesamiento, pero también disminuir las limitaciones que tienen los patólogos al implementar proyectos de análisis automático.
\endgroup
